% Abstract — DESI paper
% Target: Nature Genetics — 150-250 words
% No citations in abstract (NG style)

The interpretation of genomic diversity patterns from the Middle Pleistocene (~1.0–0.5 Ma) is contested: site frequency spectrum analyses have inferred a severe panmictic bottleneck in human ancestors at ~930 ka, while structured coalescent models favour a two-population ancestral scenario instead. Distinguishing these hypotheses requires a method that does not assume panmixia. Here we introduce DESI (Deep-time Evolutionary Structure Inference), which tests the panmixia null by comparing within-group mean pairwise coalescence times (TMRCA) across continental populations. Under a strictly panmictic model (a single undivided population throughout), all continental groups are random draws from the same genealogy and must have equal within-group TMRCA; our simulations confirm this prediction yields a near-zero ($\sim$0.2 ka) difference regardless of bottleneck parameterization. Applying DESI to phased whole-genome data from 240 individuals across five continental super-populations from the 1000 Genomes Project, we find that within-African mean TMRCA (1,229.5 ± 6.4 ka) exceeds within-East Asian TMRCA (858.8 ± 5.4 ka) by 370.7 ± 8.4 ka, with the difference consistent across all 22 autosomes (Z = 73). Coalescent simulations demonstrate that any strictly panmictic model predicts this difference to be $\sim$0.2~ka, while a Cobraa-inspired model with reduced ancestral effective population size ($N_e = 12{,}000$) predicts 370.1~ka — a 0.2\% match to the observation. The per-window TMRCA distribution shows no East Asian genomic component coalescing below 100 ka, inconsistent with a calibrated Out-of-Africa bottleneck model. The signal is unchanged in putatively neutral windows (2.6\% change), excluding background selection as a confound. These results demonstrate that the observed $\sim$370 ka genealogical depth asymmetry between African and East Asian populations is inconsistent with any strictly panmictic model and is quantitatively reproduced by demographic parameterizations motivated by the Cobraa ancestral structure hypothesis.
